% ============================================================================
% Relatório Técnico — Refatoração Flask CRUD: SQLite → PostgreSQL + Docker
% ============================================================================
\documentclass[12pt,a4paper]{article}

% --- Pacotes ---
\usepackage[utf8]{inputenc}
\usepackage[T1]{fontenc}
\usepackage[brazil]{babel}
\usepackage{geometry}
\usepackage{graphicx}
\usepackage{hyperref}
\usepackage{listings}
\usepackage{xcolor}
\usepackage{enumitem}
\usepackage{booktabs}
\usepackage{fancyhdr}
\usepackage{titlesec}

% --- Layout ---
\geometry{left=2.5cm, right=2.5cm, top=2.5cm, bottom=2.5cm}

% --- Cores personalizadas ---
\definecolor{codebg}{RGB}{245,245,245}
\definecolor{codeframe}{RGB}{200,200,200}
\definecolor{codekeyword}{RGB}{0,0,180}
\definecolor{codecomment}{RGB}{0,128,0}
\definecolor{codestring}{RGB}{163,21,21}

% --- Configuração do lstlisting ---
\lstset{
  backgroundcolor=\color{codebg},
  frame=single,
  rulecolor=\color{codeframe},
  basicstyle=\ttfamily\footnotesize,
  keywordstyle=\color{codekeyword}\bfseries,
  commentstyle=\color{codecomment}\itshape,
  stringstyle=\color{codestring},
  breaklines=true,
  breakatwhitespace=false,
  showstringspaces=false,
  tabsize=2,
  captionpos=b,
  numbers=left,
  numberstyle=\tiny\color{gray},
  numbersep=8pt,
  xleftmargin=15pt,
  framexleftmargin=15pt,
  literate=
    {á}{{\'a}}1 {é}{{\'e}}1 {í}{{\'i}}1 {ó}{{\'o}}1 {ú}{{\'u}}1
    {ã}{{\~a}}1 {õ}{{\~o}}1 {ç}{{\c{c}}}1
    {Á}{{\'A}}1 {É}{{\'E}}1 {Í}{{\'I}}1 {Ó}{{\'O}}1 {Ú}{{\'U}}1
    {Ã}{{\~A}}1 {Õ}{{\~O}}1 {Ç}{{\c{C}}}1
}

% --- Cabeçalho e rodapé ---
\pagestyle{fancy}
\fancyhf{}
\fancyhead[L]{\small Relatório Técnico — Refatoração Flask CRUD}
\fancyhead[R]{\small\thepage}
\renewcommand{\headrulewidth}{0.4pt}

% --- Hyperlinks ---
\hypersetup{
  colorlinks=true,
  linkcolor=blue!70!black,
  urlcolor=blue!70!black,
  citecolor=blue!70!black
}

% ============================================================================
\begin{document}

% --- Capa ---
\begin{titlepage}
  \centering
  \vspace*{3cm}
  {\Huge\bfseries Guia Técnico Executável e Reprodutível\par}
  \vspace{0.5cm}
  {\Large\bfseries Aluno: Kevin Grünfeld Strey\par}
  \vspace{0.5cm}
  {\LARGE Refatoração de CRUD implementado com Flask\par}
  \vspace{0.3cm}
  {\Large SQLite $\rightarrow$ PostgreSQL + Docker\par}
  \vspace{2cm}
  {\large
    \textbf{Repositório Original:}\\[4pt]
    \url{https://github.com/iamsamimsarkar/flask-crud-sqlite.git}
  \par}
  \vspace{0.65cm}
  {\large
    \textbf{Repositório Refatorado:}\\[4pt]
    \url{https://github.com/KevinStrey/flask-crud-sqlite}
  \par}
  \vspace{2cm}
  {\large \today\par}
  \vfill
\end{titlepage}

% --- Sumário ---
\tableofcontents
\newpage

% ============================================================================
% SEÇÃO 1 — Contextualização e Motivação
% ============================================================================
\section{Contextualização e Motivação}
\label{sec:contextualizacao}

\subsection{Descrição do Projeto Original}

O repositório
\href{https://github.com/iamsamimsarkar/flask-crud-sqlite.git}{\texttt{flask-crud-sqlite}}
contém uma aplicação web construída com o \textit{framework} Flask (Python) que implementa
operações CRUD (\textit{Create, Read, Update, Delete}) sobre uma entidade \texttt{User}.

A arquitetura original apresenta as seguintes características:

\begin{itemize}[nosep]
  \item \textbf{Banco de dados:} SQLite, armazenado localmente no diretório
        \texttt{instance/database.db}.
  \item \textbf{ORM:} Flask-SQLAlchemy, com modelo único (\texttt{User}) contendo os
        campos \texttt{id}, \texttt{name}, \texttt{city} e \texttt{contact}.
  \item \textbf{Servidor:} Servidor de desenvolvimento embutido do Flask
        (\texttt{app.run(debug=True)}).
  \item \textbf{Configuração:} Valores \textit{hard-coded} diretamente no arquivo
        \texttt{app.py} (string de conexão, chave secreta).
  \item \textbf{Infraestrutura:} Nenhuma — a aplicação é executada diretamente na
        máquina do desenvolvedor, sem padronização de ambiente.
\end{itemize}

\subsection{Cenário Atual vs.\ Cenário Alvo}

A tabela~\ref{tab:comparativo} resume a comparação entre o estado atual e o estado
desejado após a refatoração.

\begin{table}[h!]
  \centering
  \caption{Comparativo: estado atual vs.\ estado alvo.}
  \label{tab:comparativo}
  \begin{tabular}{lll}
    \toprule
    \textbf{Aspecto}       & \textbf{Atual}          & \textbf{Alvo} \\
    \midrule
    Banco de dados         & SQLite (arquivo local)  & PostgreSQL (contêiner) \\
    Config. do ambiente    & Hard-coded     & Variáveis de ambiente \\
    Ambiente de execução   & Máquina local           & Docker / Docker Compose \\
    Escalabilidade         & Limitada                & Horizontal (contêineres) \\
    \bottomrule
  \end{tabular}
\end{table}

\subsection{Benefícios Esperados}

\begin{enumerate}[nosep]
  \item \textbf{Escalabilidade:} O PostgreSQL suporta conexões concorrentes,
        transações ACID completas e recursos avançados de indexação, tornando a
        aplicação pronta para cenários de produção.
  \item \textbf{Ambiente padronizado:} Com Docker e Docker Compose, qualquer membro
        da equipe pode reproduzir o ambiente completo (API + banco) com um único
        comando \texttt{docker-compose up}.
  \item \textbf{Facilidade de deploy:} A conteinerização permite integração direta
        com pipelines de CI/CD e plataformas de orquestração como Kubernetes, AWS ECS
        ou Google Cloud Run.
  \item \textbf{Segurança:} Credenciais e configurações sensíveis passam a ser
        gerenciadas por variáveis de ambiente, removendo-as do código-fonte.
\end{enumerate}

% ============================================================================
% SEÇÃO 2 — Ambiente e Pré-requisitos
% ============================================================================
\section{Ambiente e Pré-requisitos}
\label{sec:prerequisitos}

\subsection{Dependências de Software}

A tabela~\ref{tab:dependencias} lista todas as ferramentas e bibliotecas necessárias
para reproduzir o ambiente refatorado.

\begin{table}[h!]
  \centering
  \caption{Mapeamento de dependências do projeto.}
  \label{tab:dependencias}
  \begin{tabular}{llp{7cm}}
    \toprule
    \textbf{Dependência}         & \textbf{Versão mínima} & \textbf{Finalidade} \\
    \midrule
    Python                       & 3.11                   & Interpretador da aplicação Flask. \\
    pip                          & 23.0                   & Gerenciador de pacotes Python. \\
    Flask                        & 3.0                    & \textit{Framework} web para a API. \\
    Flask-SQLAlchemy             & 3.1                    & ORM para abstração do banco de dados. \\
    psycopg2-binary              & 2.9                    & Driver PostgreSQL para Python/SQLAlchemy. \\
    gunicorn                     & 21.2                   & Servidor WSGI de produção. \\
    Docker Engine                & 24.0                   & \textit{Runtime} de contêineres. \\
    Docker Compose               & 2.20                   & Orquestrador de contêineres multi-serviço. \\
    Git                          & 2.40                   & Controle de versão do código-fonte. \\
    \bottomrule
  \end{tabular}
\end{table}

\subsection{Verificação do Ambiente}

Antes de iniciar a refatoração, verifique se as ferramentas estão instaladas
executando os comandos abaixo no terminal:

\begin{lstlisting}[language=bash, caption={Verificação das dependências do ambiente.}]
python --version        # Python 3.11+
pip --version           # pip 23.0+
docker --version        # Docker Engine 24.0+
docker compose version  # Docker Compose v2.20+
git --version           # Git 2.40+
\end{lstlisting}

\subsection{Estrutura Original do Projeto}

A estrutura de diretórios do repositório clonado é apresentada a seguir:

\begin{lstlisting}[caption={Árvore de diretórios do projeto original.}]
flask-crud-sqlite/
|-- app.py                  # Aplicacao Flask (ponto de entrada)
|-- instance/
|   +-- database.db         # Banco SQLite (gerado automaticamente)
|-- static/
|   +-- style.css           # Folha de estilos
|-- templates/
|   |-- index.html          # Listagem de usuarios
|   |-- add_user.html       # Formulario de cadastro
|   +-- edit_user.html      # Formulario de edicao
+-- README.md
\end{lstlisting}

% ============================================================================
% SEÇÃO 3 — Passo a Passo Executável
% ============================================================================
\section{Passo a Passo Executável}
\label{sec:passo-a-passo}

Esta seção apresenta, de forma sequencial e reprodutível, todas as modificações
realizadas no projeto. Cada subseção contém o código completo do arquivo alterado
ou criado, acompanhado de explicações técnicas.

% --------------------------------------------------------------------------
\subsection{Passo 1 — Atualização das Dependências (\texttt{requirements.txt})}

O arquivo \texttt{requirements.txt} foi criado (o projeto original não o possuía)
com as seguintes dependências:

\begin{lstlisting}[language=Python, caption={\texttt{requirements.txt} — dependências do projeto refatorado.}]
Flask==3.0.3
Flask-SQLAlchemy==3.1.1
psycopg2-binary==2.9.9
gunicorn==22.0.0
\end{lstlisting}

\textbf{Justificativas:}
\begin{itemize}[nosep]
  \item \texttt{psycopg2-binary}: driver PostgreSQL para Python, necessário para que
        o SQLAlchemy se comunique com o PostgreSQL.
  \item \texttt{gunicorn}: servidor WSGI de produção que substitui o servidor de
        desenvolvimento embutido do Flask, oferecendo melhor desempenho e
        estabilidade.
\end{itemize}

% --------------------------------------------------------------------------
\subsection{Passo 2 — Refatoração do Código Fonte (\texttt{app.py})}

O arquivo \texttt{app.py} foi modificado para ler configurações de variáveis de
ambiente, mantendo compatibilidade retroativa com SQLite como \textit{fallback}
para desenvolvimento local.

\begin{lstlisting}[language=Python, caption={\texttt{app.py} — código refatorado com suporte a variáveis de ambiente e testes.}]
import os
from flask import Flask, request, redirect, url_for, render_template
from flask_sqlalchemy import SQLAlchemy

app = Flask(__name__)

# --- Configuração via variáveis de ambiente ---
app.config['SQLALCHEMY_DATABASE_URI'] = os.environ.get(
    'DATABASE_URL',
    'sqlite:///database.db'  # fallback para desenvolvimento local sem Docker
)
app.config['SQLALCHEMY_TRACK_MODIFICATIONS'] = False
app.config['SECRET_KEY'] = os.environ.get('SECRET_KEY', 'dev-secret-key')

db = SQLAlchemy(app)


class User(db.Model):
    id = db.Column(db.Integer, primary_key=True)
    name = db.Column(db.String(100), nullable=False)
    city = db.Column(db.String(120), nullable=False)
    contact = db.Column(db.String(100), nullable=False)


with app.app_context():
    db.create_all()


@app.route("/")
def index():
    users = User.query.all()
    return render_template("index.html", users=users)


@app.route("/add", methods=['GET', 'POST'])
def add_user():
    if request.method == "POST":
        name = request.form['name']
        city = request.form['city']
        contact = request.form['contact']
        new_user = User(name=name, city=city, contact=contact)
        try:
            db.session.add(new_user)
            db.session.commit()
            return redirect(url_for("index"))
        except Exception:
            db.session.rollback()
            return 'Error adding User!'
    return render_template("add_user.html")


@app.route("/edit/<int:id>", methods=['GET', 'POST'])
def edit_user(id):
    user = db.get_or_404(User, id)
    if request.method == "POST":
        user.name = request.form['name']
        user.city = request.form['city']
        user.contact = request.form['contact']
        try:
            db.session.commit()
            return redirect(url_for('index'))
        except Exception:
            db.session.rollback()
            return 'Error editing User!'
    return render_template("edit_user.html", user=user)


@app.route("/delete/<int:id>")
def delete_user(id):
    user = db.get_or_404(User, id)
    try:
        db.session.delete(user)
        db.session.commit()
        return redirect(url_for('index'))
    except Exception:
        db.session.rollback()
        return 'Error deleting User!'


if __name__ == "__main__":
    import sys
    if len(sys.argv) > 1 and sys.argv[1] == "test":
        import unittest
        loader = unittest.TestLoader()
        suite = loader.discover('.', pattern='test_*.py')
        runner = unittest.TextTestRunner(verbosity=2)
        result = runner.run(suite)
        sys.exit(0 if result.wasSuccessful() else 1)
    app.run(debug=True, host='0.0.0.0')
\end{lstlisting}

\textbf{Alterações principais em relação ao código original:}
\begin{enumerate}[nosep]
  \item \texttt{SQLALCHEMY\_DATABASE\_URI} agora lê a variável de ambiente
        \texttt{DATABASE\_URL}. Se não estiver definida, utiliza SQLite como
        \textit{fallback}, preservando a compatibilidade local.
  \item \texttt{SECRET\_KEY} também é externalizada via variável de ambiente.
  \item Adicionado \texttt{db.session.rollback()} nos blocos \texttt{except} de todas
        as rotas de mutação de dados (\texttt{add}, \texttt{edit} e \texttt{delete}) para
        garantir integridade transacional em caso de erro.
  \item \texttt{User.query.get()} substituído por \texttt{db.get\_or\_404()}
        para retornar HTTP 404 de forma padronizada e segura quando o registro não existir.
  \item Inclusão de suporte a execução de testes de regressão via CLI (\texttt{python app.py test})
        no bloco principal \texttt{\_\_main\_\_}.
  \item \texttt{host='0.0.0.0'} adicionado ao \texttt{app.run()} para permitir
        conexões de fora do contêiner Docker.
\end{enumerate}

% --------------------------------------------------------------------------
\subsection{Passo 3 — Criação do Dockerfile}

O \texttt{Dockerfile} utiliza a técnica de \textit{multi-stage build} para produzir
uma imagem final enxuta e segura.

\begin{lstlisting}[language=bash, caption={\texttt{Dockerfile} — build multi-estágio otimizado.}]
# ---- Estágio 1: Construção ----
FROM python:3.11-slim AS builder

WORKDIR /app

# Instala dependências do sistema para compilar psycopg2
RUN apt-get update && \
    apt-get install -y --no-install-recommends gcc libpq-dev && \
    rm -rf /var/lib/apt/lists/*

COPY requirements.txt .
RUN pip install --no-cache-dir --prefix=/install -r requirements.txt

# ---- Estágio 2: Produção ----
FROM python:3.11-slim

WORKDIR /app

# Instala apenas a libpq (runtime)
RUN apt-get update && \
    apt-get install -y --no-install-recommends libpq5 && \
    rm -rf /var/lib/apt/lists/*

# Copia as dependências instaladas do estágio anterior
COPY --from=builder /install /usr/local

# Copia o código da aplicação
COPY . .

# Cria um usuário não-root para segurança
RUN useradd --create-home appuser
USER appuser

EXPOSE 5000

# Inicia com Gunicorn (servidor WSGI de produção)
CMD ["gunicorn", "--bind", "0.0.0.0:5000", "--workers", "3", "app:app"]
\end{lstlisting}

\textbf{Decisões de projeto:}
\begin{itemize}[nosep]
  \item \textbf{Multi-stage build:} O estágio \texttt{builder} contém compiladores
        (\texttt{gcc}, \texttt{libpq-dev}) necessários apenas durante a instalação
        do \texttt{psycopg2-binary}. A imagem final não carrega esses pacotes,
        resultando em uma imagem significativamente menor.
  \item \textbf{Usuário não-root:} O contêiner executa sob o usuário
        \texttt{appuser}, seguindo o princípio de menor privilégio.
  \item \textbf{Gunicorn com 3 workers:} Configuração adequada para ambientes de
        desenvolvimento e pequena produção. Em produção real, o número de
        \textit{workers} deve ser ajustado conforme os recursos disponíveis
        (recomendação: \texttt{2 * num\_cpus + 1}).
\end{itemize}

% --------------------------------------------------------------------------
\subsection{Passo 4 — Criação do Docker Compose (\texttt{docker-compose.yml})}

O arquivo \texttt{docker-compose.yml} orquestra dois serviços: o banco PostgreSQL
e a aplicação Flask.

\begin{lstlisting}[language=bash, caption={\texttt{docker-compose.yml} — orquestração dos serviços.}]
version: "3.8"

services:
  # ---- Serviço: Banco de dados PostgreSQL ----
  db:
    image: postgres:16-alpine
    container_name: flask_crud_db
    restart: unless-stopped
    environment:
      POSTGRES_USER: flask_user
      POSTGRES_PASSWORD: flask_pass
      POSTGRES_DB: flask_crud
    ports:
      - "5432:5432"
    volumes:
      - pgdata:/var/lib/postgresql/data
    healthcheck:
      test: ["CMD-SHELL", "pg_isready -U flask_user -d flask_crud"]
      interval: 5s
      timeout: 5s
      retries: 5

  # ---- Serviço: Aplicação Flask ----
  web:
    build: .
    container_name: flask_crud_web
    restart: unless-stopped
    ports:
      - "5000:5000"
    environment:
      DATABASE_URL: postgresql://flask_user:flask_pass@db:5432/flask_crud
      SECRET_KEY: super-secret-key-change-in-production
    depends_on:
      db:
        condition: service_healthy

volumes:
  pgdata:
\end{lstlisting}

\textbf{Pontos importantes:}
\begin{itemize}[nosep]
  \item \textbf{Healthcheck no PostgreSQL:} O serviço \texttt{web} só inicia após o
        PostgreSQL reportar que está pronto para aceitar conexões, evitando erros de
        conexão durante a inicialização.
  \item \textbf{Volume nomeado (\texttt{pgdata}):} Os dados do PostgreSQL persistem
        entre reinicializações dos contêineres.
  \item \textbf{Rede interna:} O Docker Compose cria automaticamente uma rede
        \textit{bridge} entre os serviços. O serviço \texttt{web} referencia o banco
        pelo hostname \texttt{db} na string de conexão.
\end{itemize}

% --------------------------------------------------------------------------
\subsection{Passo 5 — Arquivo \texttt{.dockerignore}}

Para otimizar o contexto de build, foi criado um \texttt{.dockerignore}:

\begin{lstlisting}[caption={\texttt{.dockerignore} — arquivos excluídos do build.}]
__pycache__/
*.pyc
*.pyo
instance/
.env
*.db
.git/
.gitignore
*.md
Relatorio.tex
\end{lstlisting}

% --------------------------------------------------------------------------
\subsection{Passo 6 — Comandos de Operação}

\subsubsection{Subir a aplicação}

\begin{lstlisting}[language=bash, caption={Comandos para construir e iniciar os contêineres.}]
# Construir as imagens e subir em modo detached
docker compose up --build -d

# Verificar se os contêineres estão rodando
docker compose ps

# Acompanhar os logs em tempo real
docker compose logs -f
\end{lstlisting}

Após a execução, a aplicação estará disponível em \texttt{http://localhost:5000}.

\subsubsection{Derrubar a aplicação}

\begin{lstlisting}[language=bash, caption={Comandos para parar e remover os contêineres.}]
# Parar os contêineres (preserva dados do volume)
docker compose down

# Parar e REMOVER os volumes (apaga dados do PostgreSQL)
docker compose down -v
\end{lstlisting}

\subsubsection{Estrutura Final do Projeto}

Após todas as modificações, a árvore de diretórios do projeto fica:

\begin{lstlisting}[caption={Árvore de diretórios do projeto refatorado.}]
flask-crud-sqlite/
|-- app.py                  # Aplicacao Flask (refatorado)
|-- test_app.py             # Testes de regressao automatizados (novo)
|-- requirements.txt        # Dependências Python (novo)
|-- Dockerfile              # Build da imagem Docker (novo)
|-- docker-compose.yml      # Orquestração dos servicos (novo)
|-- .dockerignore           # Exclusões do build Docker (novo)
|-- static/
|   +-- style.css           # Folha de estilos
|-- templates/
|   |-- index.html          # Listagem de usuarios
|   |-- add_user.html       # Formulario de cadastro
|   +-- edit_user.html      # Formulario de edicao
+-- README.md
\end{lstlisting}

% ============================================================================
% SEÇÃO 4 — Plano de Rollback e Testes de Validação
% ============================================================================
\section{Plano de Rollback e Testes de Validação}
\label{sec:rollback}

\subsection{Testes de Validação}

Após subir a aplicação com \texttt{docker compose up --build -d}, execute os
seguintes testes para validar que tudo funciona corretamente.

\subsubsection{Teste 1 — Verificar status dos contêineres}

\begin{lstlisting}[language=bash, caption={Verificação do estado dos contêineres.}]
docker compose ps
\end{lstlisting}

\textbf{Resultado esperado:} Ambos os serviços (\texttt{flask\_crud\_db} e
\texttt{flask\_crud\_web}) devem aparecer com status \texttt{Up} e, no caso do
banco, \texttt{(healthy)}.

\subsubsection{Teste 2 — Acessar a página inicial via navegador}

Abra \texttt{http://localhost:5000} no navegador. A página de listagem de usuários
deve ser exibida corretamente, sem erros.

% \subsubsection{Teste 3 — Testar o CRUD via linha de comando}

% \begin{lstlisting}[language=bash, caption={Testes CRUD via \texttt{curl}.}]
% # Teste de leitura (GET na página inicial)
% curl -s -o /dev/null -w "%{http_code}" http://localhost:5000/
% # Esperado: 200

% # Teste de criação (POST para adicionar usuário)
% curl -X POST http://localhost:5000/add \
%   -d "name=Teste&city=Brasilia&contact=teste@email.com" \
%   -w "\n%{http_code}"
% # Esperado: 302 (redirect após criação bem-sucedida)

% # Teste de leitura após criação
% curl -s http://localhost:5000/ | grep "Teste"
% # Esperado: o nome "Teste" deve aparecer no HTML retornado
% \end{lstlisting}

\subsubsection{Teste 3 — Verificar conexão com o PostgreSQL}

\begin{lstlisting}[language=bash, caption={Verificação direta no banco PostgreSQL.}]
docker compose exec db psql -U flask_user -d flask_crud -c "\dt"
\end{lstlisting}

\textbf{Resultado esperado:} A tabela \texttt{user} deve aparecer listada,
confirmando que o SQLAlchemy criou o esquema corretamente no PostgreSQL.

\subsubsection{Teste 4 — Verificar logs da aplicação}

\begin{lstlisting}[language=bash, caption={Inspeção de logs para identificar erros.}]
docker compose logs web
\end{lstlisting}

\textbf{Resultado esperado:} Os logs do Gunicorn devem indicar que os
\textit{workers} foram inicializados com sucesso, sem mensagens de erro de conexão
com o banco.

% --------------------------------------------------------------------------
\subsection{Plano de Rollback}

Caso a refatoração apresente problemas ou seja necessário reverter para o estado
original, siga os passos abaixo:

\subsubsection{Passo 1 — Derrubar a infraestrutura Docker}

\begin{lstlisting}[language=bash, caption={Remoção completa dos contêineres e volumes.}]
docker compose down -v
\end{lstlisting}

Este comando para todos os contêineres, remove-os e apaga o volume
\texttt{pgdata} com os dados do PostgreSQL.

\subsubsection{Passo 2 — Reverter o código-fonte via Git para o commit original}

\begin{lstlisting}[language=bash, caption={Rollback do código via Git com especificação do commit original.}]
# Reverter o repositório diretamente para o commit original específico (pré-refatoração)
git checkout aa5bafc -- .

# Remover arquivos novos não rastreados criados na refatoração (Dockerfile, etc.)
git clean -fd
\end{lstlisting}

\textbf{Detalhamento técnico da melhoria:} A indicação explícita do \textit{commit hash}
original (\texttt{aa5bafc}) substitui o uso de comandos genéricos como \texttt{git revert}
desacompanhado de parâmetros. Em históricos com múltiplos \textit{commits}, um
\texttt{git revert} genérico pode tentar desfazer apenas o último \textit{commit},
gerar conflitos de mesclagem complexos ou deixar o repositório em um estado intermediário
quebrado. O apontamento determinístico para o \textit{commit} original garante a
reversão atômica e integral da base de código.

\subsubsection{Passo 3 — Verificar o estado do repositório}

\begin{lstlisting}[language=bash, caption={Confirmação do estado limpo do repositório.}]
git status
\end{lstlisting}

\textbf{Resultado esperado:} O Git deve reportar \texttt{nothing to commit,
working tree clean}, confirmando que o repositório voltou ao estado original.

\subsubsection{Passo 4 — Testes de Regressão e Validação da Aplicação Original}

\begin{lstlisting}[language=bash, caption={Execução dos testes de regressão automatizados e inicialização original.}]
# 1. Instalar dependências da aplicação original
pip install flask flask-sqlalchemy

# 2. Executar suíte de testes de regressão pós-rollback
python -m unittest test_app.py
# ou alternativamente:
python app.py test

# 3. Iniciar a aplicação original para verificação final em desenvolvimento
python app.py
\end{lstlisting}

\textbf{Detalhamento técnico da melhoria:} Anteriormente, a validação pós-rollback
resumia-se a ``testes de sistema'' manuais (abrir o navegador e interagir com formulários),
uma metodologia imprecisa, lenta e sujeita ao esquecimento de cenários críticos de borda.
A introdução dos testes de regressão automatizados (\texttt{test\_app.py}) permite aferir
imediatamente e com 100\% de precisão que todas as quatro operações do CRUD, redirecionamentos
e persistência em SQLite continuam íntegros após o desfazimento da migração. O acesso visual
em \texttt{http://localhost:5000} torna-se uma validação complementar.

% --------------------------------------------------------------------------
\subsection{Resumo do Fluxo de Rollback}

\begin{table}[h]
  \centering
  \caption{Resumo dos comandos de rollback.}
  \vspace{0.4em}
  \label{tab:rollback}
  \begin{tabular}{clp{7.5cm}}
    \toprule
    \textbf{Ordem} & \textbf{Comando} & \textbf{Efeito} \\
    \midrule
    1 & \texttt{docker compose down -v} & Remove contêineres, redes e volume pgdata. \\
    2 & \texttt{git checkout aa5bafc -- .} & Reverte o código para o commit original estável. \\
    3 & \texttt{git clean -fd}          & Remove arquivos não rastreados criados na refatoração. \\
    4 & \texttt{git status}             & Confirma repositório limpo. \\
    5 & \texttt{python -m unittest test\_app.py} & Executa testes de regressão automatizados no SQLite. \\
    6 & \texttt{python app.py}          & Inicia a aplicação original com SQLite. \\
    \bottomrule
  \end{tabular}
\end{table}

% ============================================================================
% APÊNDICE — Quadro de Melhorias (Pós-Peer Review)
% ============================================================================
\newpage
\section{Apêndice: Quadro de Melhorias (Pós-Peer Review)}
\label{sec:melhorias-peer-review}

\subsection{Contextualização e Síntese dos Feedbacks Recebidos}

Após a elaboração da primeira versão deste relatório técnico de refatoração, o documento
foi submetido à dinâmica de avaliação por pares (\textit{peer review}) entre os colegas
de turma. A avaliação recebida foi amplamente positiva, ressaltando o cumprimento rigoroso
de todos os critérios estabelecidos no enunciado e a clareza didática da apresentação:

\begin{quote}
\textit{``Trabalho muito bem detalhado e muito bem organizado. Cumpre todos os
requisitos propostos e de maneira fácil de ser compreendido. Plano muito bem
detalhado e compreensível. Os snippets de código são muito bem organizados, tudo é
apresentado numa sequência clara e os items estão todos presentes. Ficou muito bom.
Ele atende todos os tópicos solicitados no enunciado, desde a contextualização e os
pré-requisitos até o passo a passo do roteiro de migração, o plano de rollback e os
testes de validação.''}
\end{quote}

\begin{quote}
\textit{``Gostei bastante da organização do documento LaTeX. As tabelas comparando o
antes e depois, os blocos de código e as explicações deixam o contéudo bem fácil de
compreender. No geral, o relatório ficou bem completo, claro e estruturado para
orientar uma possível execução da migração. Parabéns pelo empenho!''}
\end{quote}

\subsection{Identificação de Oportunidades e Decisões de Complementação}

Não obstante o retorno positivo unânime e a ausência de apontamentos de falhas estruturais,
o próprio processo de revisão dos relatórios elaborados pelos colegas de turma forneceu
valiosos \textit{insights} comparativos. A partir dessa análise cruzada de metodologias,
decidiu-se enriquecer o projeto com duas complementações substanciais:

\begin{enumerate}
  \item \textbf{Criação de casos de teste automatizados para regressão pós-rollback:}
        Antes da revisão, a validação da aplicação após o procedimento de rollback era
        conduzida apenas mediante testes de sistema manuais via navegador web. Essa
        prática apresentava baixa precisão, dependia de inspeção visual subjetiva e não
        oferecia reprodutibilidade técnica. A partir do \textit{insight} adquirido nos
        relatórios dos colegas, foi implementada uma suíte formal de testes de regressão
        (\texttt{test\_app.py}), permitindo atestar com rigor que a aplicação original
        preserva sua integridade funcional e operacional com SQLite.
  \item \textbf{Adição de commit específico no processo de rollback via Git:}
        Na abordagem preliminar, a reversão de versão no Git era descrita de forma genérica
        com \texttt{git revert} (ou comandos de descarte de modificações sem fixação
        de \textit{hash}). No entanto, em históricos de desenvolvimento com múltiplos
        \textit{commits}, o \texttt{git revert} pode gerar comportamentos inesperados,
        conflitos de mesclagem ou reversões parciais caso não receba exatamente quais
        \textit{commits} devem ser revertidos. Para eliminar essa margem de incerteza,
        foi adicionado o identificador do \textit{commit} original pré-refatoração
        (\texttt{aa5bafc}), tornando o procedimento 100\% determinístico, reproduzível
        e seguro.
\end{enumerate}

\subsection{Quadro Resumo de Melhorias (Pós-Peer Review)}

A tabela~\ref{tab:quadro-melhorias} sintetiza as complementações realizadas no código
e no relatório técnico em decorrência da avaliação por pares.

\begin{table}[h!]
  \centering
  \caption{Quadro de Melhorias implementadas pós-Peer Review.}
  \vspace{0.4em}
  \label{tab:quadro-melhorias}
  \small
  \begin{tabular}{p{3.0cm}p{3.2cm}p{3.8cm}p{4.2cm}}
    \toprule
    \textbf{Componente / Área} & \textbf{Abordagem Anterior} & \textbf{Insight do Peer Review} & \textbf{Melhoria Implementada} \\
    \midrule
    \textbf{Validação Pós-Rollback} &
    Teste de sistema manual no navegador (\texttt{http://localhost:5000}). &
    Testes manuais são imprecisos, subjetivos e sujeitos a falhas humanas não detectadas. &
    Criação do arquivo \texttt{test\_app.py} com 7 testes de regressão automatizados cobrindo todo o ciclo CRUD. \\
    \addlinespace
    \textbf{Controle de Versão (Git Rollback)} &
    Comandos genéricos (\texttt{git revert} / checkout sem hash explícito). &
    Comandos genéricos podem reverter commits indesejados ou gerar conflitos de merge inesperados. &
    Fixação do commit hash original (\texttt{git checkout aa5bafc -- .}), garantindo rollback determinístico. \\
    \addlinespace
    \textbf{Tratamento Transacional (\texttt{app.py})} &
    \texttt{db.session.rollback()} presente apenas na criação (\texttt{/add}). &
    Possibilidade de estados inconsistentes na sessão em caso de falhas durante testes de regressão. &
    Extensão do \texttt{try/except} com rollback para as rotas \texttt{/edit} e \texttt{/delete}, e uso de \texttt{db.get\_or\_404}. \\
    \addlinespace
    \textbf{Execução de Testes via CLI} &
    Inexistente (execução apenas via servidor web em execução). &
    Necessidade de execução rápida e automatizável de testes diretamente pelo ponto de entrada. &
    Implementação do subcomando \texttt{python app.py test} no bloco \texttt{\_\_main\_\_} de \texttt{app.py}. \\
    \bottomrule
  \end{tabular}
\end{table}

\subsection{Detalhamento das Melhorias no Código e Procedimentos}

\subsubsection{1. Suíte de Testes de Regressão (\texttt{test\_app.py})}

A suíte de testes de regressão foi desenvolvida sobre o módulo padrão \texttt{unittest}
do Python, garantindo portabilidade sem necessidade de bibliotecas externas complexas.
O banco de dados é isolado durante os testes via variável \texttt{DATABASE\_URL},
assegurando que o arquivo de produção/desenvolvimento \texttt{instance/database.db}
permaneça intacto. Foram contemplados 7 casos de teste fundamentais:

\begin{itemize}[nosep]
  \item \texttt{test\_01\_index\_empty}: Confirma HTTP 200 na listagem inicial sem registros.
  \item \texttt{test\_02\_add\_user\_success}: Valida criação de usuário via POST em \texttt{/add} e persistência no banco.
  \item \texttt{test\_03\_edit\_user\_view}: Valida pré-carregamento correto dos dados na tela de edição.
  \item \texttt{test\_04\_edit\_user\_post\_success}: Valida persistência das alterações cadastrais.
  \item \texttt{test\_05\_delete\_user\_success}: Confirma remoção do usuário do banco de dados e redirecionamento.
  \item \texttt{test\_06\_multiple\_users\_persistence}: Valida inserção, contagem e persistência simultânea de múltiplos usuários.
  \item \texttt{test\_07\_full\_regression\_crud\_cycle}: Executa o ciclo completo de ponta a ponta (Create $\rightarrow$ Read $\rightarrow$ Update $\rightarrow$ Delete).
\end{itemize}

Os testes podem ser invocados com:
\begin{lstlisting}[language=bash, caption={Comandos para disparo dos testes de regressão.}]
python -m unittest test_app.py
# ou diretamente via app.py:
python app.py test
\end{lstlisting}

\subsubsection{2. Rollback com Commit Hash Específico}

A especificação do \textit{commit} \texttt{aa5bafc} no Passo 2 do Plano de Rollback
assegura que qualquer intervenção de reversão restaure pontualmente o repositório
ao estado exato entregue pelo projeto base, neutralizando qualquer deriva de código
decorrente de múltiplos \textit{commits} intermediários realizados pela equipe.

% ============================================================================
\end{document}
